\section{Introduction}
\label{sec:introduction}

Medical image computing methods are commonly developed from heterogeneous
clinical data collected across patients, scanners, sites, and time.  These
sources of variation make evaluation design inseparable from model design: a
high score is not evidence of clinical utility when related observations occur
in both training and test sets, when model selection uses the test set, or when
important patient subgroups are absent from the evaluation.  Strong
segmentation architectures such as U-Net~\cite{ronneberger2015unet}, 3D
U-Net~\cite{cicek20163d}, and nnU-Net~\cite{isensee2021nnu} illustrate the value
of carefully matched architectures and training pipelines, but fair comparison
still depends on a shared and transparent protocol.

A clear MICCAI paper should answer three questions early.  First, what
clinically or scientifically meaningful problem is unresolved?  Second, why do
existing methods fail to resolve it?  Third, which evidence would distinguish
the proposed explanation or method from plausible alternatives?  The final
paragraph of the introduction should state contributions as testable claims,
not as a list of implementation components.  A suitable pattern is:

\begin{itemize}
  \item a methodological contribution linked to a specific limitation of prior
        work;
  \item an evaluation design that tests generalization under clinically
        relevant variation; and
  \item an analysis that explains where the method succeeds, where it fails,
        and how certain those conclusions are.
\end{itemize}

The remainder of the paper should then provide evidence for these claims in
the same order.  Related work establishes the gap, the method defines the
intervention, and the experiments isolate its effect.
